\documentclass[10pt,a4paper]{article}
\usepackage[T1]{fontenc}\usepackage{lmodern,amsmath,amssymb,amsthm,mathtools,microtype}
\usepackage[margin=24mm]{geometry}\usepackage[colorlinks=true,linkcolor=blue,urlcolor=blue,citecolor=blue]{hyperref}
\newtheorem{theorem}{Theorem}[section]\newtheorem{lemma}[theorem]{Lemma}
\newcommand{\Q}{\mathbb Q}\newcommand{\Z}{\mathbb Z}
\setlength{\emergencystretch}{2em}
\title{Pole equations from simultaneous congruent-number twists}
\author{}\date{14 November 2026}
\begin{document}\maketitle
\begin{abstract}
Three equations forbid a negative odd valuation of their rational input
at every odd prime. For prepared positive inputs less than one, the
simultaneous-twist theorem of Jaswal supplies rational witnesses.
The application checks the precise square classes, local congruences
and scaling of both elliptic points, and treats the square input
separately. Seven auxiliary rational coordinates give a quadratic
presentation, with smaller coefficients than the previous pole system.
The pole replacement does not reduce the coordinate count on its own.
\end{abstract}

\section{The equations and their local implication}
For $b\in\Q$, use rational $e,p,u,v,q$ and require
\begin{equation}\label{eq:pole}
 e^2+bp^2=1,\qquad uv=e(1-b),\qquad u^2-v^2=bq^2.
\end{equation}
No nonzero restriction is imposed for the local implication.

\begin{theorem}\label{thm:local}
If \eqref{eq:pole} has a rational solution and $b\ne0$, then
$v_\ell(b)$ is not a negative odd integer for any odd prime $\ell$.
\end{theorem}
\begin{proof}
Suppose $\nu=v_\ell(b)<0$ is odd. The first equation has terms of
different valuation parity whenever both are nonzero. It forces
$v_\ell(e)=0$ and $v_\ell(bp^2)>0$ (with the latter interpreted as
infinity when $p=0$): a negative smaller valuation could not cancel,
and a zero valuation on $bp^2$ is impossible by parity.
Thus $e\ne0$ and $v_\ell(uv)=v_\ell(e(1-b))=\nu$, so both $u,v$
are nonzero and their valuations have odd sum. In particular they are
distinct. The valuation of $u^2-v^2$ is then the smaller of two distinct
even integers, and is even. The right side $bq^2$ is either zero or has
odd valuation. Both possibilities contradict the nonzero even-valuation
left side.
\end{proof}

\section{The exact simultaneous-twist input}
Write $E_r:y^2=x^3-r^2x$. We use the following statement, exactly the
scope of Jaswal's Theorem~1.2 in Hexagon version 2~\cite{Jaswal}:
for coprime positive squarefree integers $\beta>1$, $\theta\ge1$,
with $\beta\equiv\theta\equiv1\pmod8$, each a square modulo every
prime dividing the other, there is a positive squarefree integer $r$
such that $E_r,E_{r\theta}$ both have an infinite-order rational point
whose abscissa lies in $\beta\Q^{\times2}$. The source gives analytic
and algebraic rank one as well, but those stronger conclusions are not
needed for the following application.

We retain this stated theorem as the input to the realization result.
The application below is proved in full; it is not a claim that every
part of the source's external arithmetic proof has been independently
re-audited. No use is made of its separate contact-height or
prime-pattern statements. In particular, the theorem does not assert
that arbitrary different prescribed square classes share witnesses.

\begin{theorem}\label{thm:complete}
Assume the simultaneous-twist statement above. Suppose $b\in\Q$
has square reduced denominator and satisfies
\[
 0<b<1,\qquad b\in\Q_2^{\times2},\qquad v_2(b)\ge3.
\]
Then \eqref{eq:pole} has a rational solution with $upq\ne0$.
\end{theorem}
\begin{proof}
Write $b=\beta z^2$ and $1-b=\theta c^2$, where $\beta,\theta$
are positive squarefree integers and $z,c\in\Q^\times$. More
explicitly, if $b=A/B^2$ in lowest terms, then $\beta$ is the squarefree
part of $A$, and $\theta$ that of $B^2-A$. They are coprime: a prime
dividing both would divide $A$ and $B$, which is impossible.
At each odd prime dividing $\beta$, the value $1-b$ has square unit
residue, so $\theta$ is a square modulo that prime. Conversely at a
prime dividing $\theta$, $b$ has square unit residue and $\beta$ is a
square. These primes cannot divide the reduced denominator.

Because $b$ is a 2-adic square, $v_2(b)$ is even and is at least four.
The square denominator is therefore odd in reduced form. The squarefree
part $\beta$ is odd and is $1\pmod8$, as follows from the square unit
part of $b$. Also $1-b\equiv1\pmod {16}$ in $\Z_2$, so its odd
squarefree part $\theta$ is $1\pmod8$.

If $\beta=1$, put $e=0$, $p=1/z$, $u=z$, $v=0$ and $q=1$.
This proves the claim directly and does not apply a theorem requiring
$\beta>1$. Otherwise apply the simultaneous-twist input to
$\beta,\theta$. Scale the point on $E_{r\theta}$ by
$(x,y)\mapsto(c^2x,c^3y)$, obtaining a point on $E_{r(1-b)}$.
Both abscissas remain in $\beta\Q^{\times2}=b\Q^{\times2}$.
Write the resulting points as $(ba_1^2,y_1)$ on $E_r$ and
$(ba_2^2,y_2)$ on $E_{r(1-b)}$, with $a_1a_2y_1y_2\ne0$.
Their ordinates are nonzero because the points have infinite order.
Set
\[
 e=\frac{r}{ba_1^2},\quad p=\frac{y_1}{b^2a_1^3},\quad
 u=\frac{a_2}{a_1},\quad v=\frac{e(1-b)}u,\quad
 q=\frac{y_2}{b^2a_1^3u^2}.
\]
The first curve equation, divided by $b^3a_1^6$, gives
$bp^2=1-e^2$. The second, divided by $b^3a_1^6u^2$, gives
$bq^2=u^2-e^2(1-b)^2/u^2=u^2-v^2$. The middle equation is the
chosen definition of $v$. All divisions are by quantities just shown
nonzero, and $upq\ne0$.
\end{proof}

\section{Quadratic presentation and coefficient comparison}
Retain two additional variables $p_2,q_2$ and impose five quadrics:
\begin{equation}\label{eq:quadrics}
 p_2=p^2,\quad q_2=q^2,\quad e^2+bp_2=1,\quad
 uv=e(1-b),\quad u^2-v^2=bq_2.
\end{equation}
There are seven auxiliary coordinates, the same number as before.
The replacement changes three rows per detector while keeping the two
square definitions. It uses no fixed factor 17, and every displayed
coefficient has absolute value one. The new prepared range is $(0,1)$;
in the abc quartic the actual pole input tends to zero on the selected
real neighborhoods, so this range can be imposed before prime selection.
The old extra squareclass conditions at 3 and 17 may be retained but
are not consumed by Theorem~\ref{thm:complete}.

At a prime with positive odd contact order, the unchanged trace equations
force a nonzero pole input of odd valuation. Theorem~\ref{thm:local}
therefore gives the same strict local inequalities as the older pole
system. No assertion about a zero pole input is needed away from these
consumed primes. The new global zero set is not claimed identical to
the old pole equations for every $b$; soundness and completeness of the
whole construction are proved separately.

In the unchanged 81-coordinate reference presentation the full quartic
has exact height $\max(4153456,B^2+630511)$. In the compressed
79-coordinate presentation the reference equations, not the pole
coefficients, dominate the height. Both statements are verified after
expanding and collecting all coefficients in their respective complete
tables. They do not give an effective sufficient defining shift.

\begin{thebibliography}{9}
\bibitem{Jaswal} M. Jaswal, \emph{A proof that $\Z$ is Diophantine in $\Q$
through congruent-number twists}, Hexagon \texttt{2611.00331v2},
12 November 2026, Theorem~1.2, printed p.~4.
\url{https://hexagonmath.org/pdf/2611.00331v2}.
\bibitem{Quartic} \emph{A 79-variable quartic under rational $abc$},
project companion, \nolinkurl{quartic_abc79_20261114_v27.tex}.
\end{thebibliography}
\end{document}
